% Part: first-order-logic
% Chapter: completeness
% Section: introduction

\documentclass[../../../include/open-logic-section]{subfiles}

\begin{document}

\iftag{FOL}
      {\olfileid{fol}{com}{int}}
      {\olfileid{pl}{com}{int}}

\olsection{ভূমিকা}

পূর্ণতা উপপাদ্য যুক্তিবিদ্যার সবচেয়ে মৌলিক ফলগুলির একটি। এর দুটি রূপ
আছে, এবং আমরা প্রমাণ করব যে তারা সমতুল্য। প্রথম রূপে এটি অর্থগত
অনুসিদ্ধান্ত ও আমাদের !!{derivation} পদ্ধতির সম্পর্ক সম্বন্ধে একটি মৌলিক
কথা বলে: কোনো !!a{sentence}~$!A$ যদি কয়েকটি !!{sentence}s~$\Gamma$ থেকে
অনুসৃত হয়, তবে $\Gamma \Proves !A$ প্রতিষ্ঠা করে এমন একটি
!!a{derivation}-ও আছে। অতএব, প্রকৃতপক্ষে অনুসৃত নয় এমন কিছু প্রমাণ না
করেও !!{derivation} পদ্ধতিটি যতটা শক্তিশালী হওয়া সম্ভব, ঠিক ততটাই
শক্তিশালী।

দ্বিতীয় রূপে একে মডেলের অস্তিত্ববিষয়ক ফল হিসেবে বলা যায়: !!{sentence}s-এর
প্রতিটি সঙ্গত সেট পরিতৃপ্তিযোগ্য। সঙ্গতি একটি প্রমাণতাত্ত্বিক ধারণা: এর
অর্থ, আমাদের !!{derivation} পদ্ধতি বিশেষ ধরনের কিছু !!{derivation}s
উৎপন্ন করতে অক্ষম। কিন্তু~$\Gamma$ থেকে বিশেষ ধরনের কোনো
!!{derivation}s নেই বলেই যে এমন একটি
\iftag{FOL}{!!a{structure}~$\Struct{M}$}{!!{valuation}{~$\pAssign{v}$}}
থাকবে যাতে $\iftag{FOL}{\Sat{M}}{\pSat{v}}{\Gamma}$, তার নিশ্চয়তা কে
দেবে? পূর্ণতা উপপাদ্য প্রথম প্রমাণিত হওয়ার আগে—বস্তুত, এখনকার
!!{derivation} পদ্ধতিগুলি গড়ে ওঠারও আগে—মহান জার্মান গণিতবিদ ডেভিড
হিলবার্টের মত ছিল যে গাণিতিক তত্ত্বগুলির সঙ্গতি তাদের আলোচ্য বস্তুগুলির
অস্তিত্বের নিশ্চয়তা দেয়। গটলব ফ্রেগেকে লেখা এক চিঠিতে তিনি কথাটি এভাবে
বলেছিলেন:
\begin{quote}
  যথেচ্ছভাবে দেওয়া স্বতঃসিদ্ধগুলি তাদের সব অনুসিদ্ধান্তসহ যদি পরস্পরের
  বিরোধী না হয়, তবে সেগুলি সত্য এবং স্বতঃসিদ্ধগুলি যে বস্তুগুলিকে
  সংজ্ঞায়িত করে সেগুলির অস্তিত্ব আছে। আমার কাছে এটিই সত্যতা ও অস্তিত্বের
  মানদণ্ড।
\end{quote}
ফ্রেগে এর তীব্র বিরোধিতা করেছিলেন। পূর্ণতা উপপাদ্যের দ্বিতীয় রূপটি অন্তত
এই অর্থে হিলবার্টকে সঠিক প্রমাণ করে যে স্বতঃসিদ্ধগুলি সঙ্গত হলে এমন
\emph{কোনো-না-কোনো} \iftag{FOL}{!!{structure}}{!!{valuation}} আছে যা
সবগুলিকে সত্য করে।

পূর্ণতা উপপাদ্য—অথবা আরও ঠিকভাবে বললে, তার প্রমাণ—শুধু এই কারণগুলির জন্যই
গুরুত্বপূর্ণ নয়। এর কয়েকটি গুরুত্বপূর্ণ পরিণাম আছে; তাদের কিছু আমরা
আলাদাভাবে আলোচনা করব। যেমন, $\Gamma \Proves !A$ দেখায় এমন প্রত্যেক
!!{derivation} সসীম, তাই সেটি~$\Gamma$-র কেবল সসীমসংখ্যক !!{sentence}s
ব্যবহার করতে পারে। ফলে পূর্ণতা উপপাদ্য থেকে পাই: $!A$ যদি~$\Gamma$-র
অনুসিদ্ধান্ত হয়, তবে সেটি ইতিমধ্যেই~$\Gamma$-র কোনো সসীম উপসেটের
অনুসিদ্ধান্ত। একে \emph{সংহতি} বলা হয়। সমতুল্যভাবে, $\Gamma$-র প্রতিটি
সসীম উপসেট সঙ্গত হলে $\Gamma$ নিজেও সঙ্গত।

!!{derivation}s-এর পথ ঘুরে সংহতি উপপাদ্য পূর্ণতা উপপাদ্য থেকে অনুসৃত হলেও,
পূর্ণতা উপপাদ্যের \emph{প্রমাণটি} ব্যবহার করে একে সরাসরিও প্রতিষ্ঠা করা
যায়। কারণ সেই প্রমাণটি বিশেষ একটি ধর্ম—সঙ্গতি—বিশিষ্ট !!{sentence}s-এর
একটি সেট নিয়ে, সেই সেট থেকে বিশেষ কিছু ধর্মবিশিষ্ট একটি !!a{structure}
নির্মাণ করে (এক্ষেত্রে গঠনটি সেটটিকে পরিতৃপ্ত করে)। সঙ্গত সেটের বদলে
“সসীমভাবে পরিতৃপ্তিযোগ্য” !!{sentence}s-এর সেট থেকে শুরু করলে প্রায় একই
নির্মাণ দিয়ে সংহতি সরাসরি প্রতিষ্ঠা করা যায়।
\iftag{FOL}{এই নির্মাণ থেকে আরও ফল পাওয়া যায়; যেমন, !!{sentence}s-এর
  যেকোনো পরিতৃপ্তিযোগ্য সেটের একটি সসীম অথবা !!{denumerable}s মডেল থাকে।
  (এই ফলটিকে লোয়েনহাইম--স্কোলেম উপপাদ্য বলা হয়।) সাধারণভাবে,
  !!{sentence}s-এর সেট থেকে !!{structure}s নির্মাণ যুক্তিবিদ্যায়, এমনকি
  কখনো কখনো দর্শনেও, প্রায়ই ব্যবহৃত হয়।}{}

\end{document}
